% =====================================================================
%  LDA-T318 — Lecture 1 — Language models
%  Tuesday 27 October 2026, 16:15–17:45, Metsätalo B525
%  Build:  latexmk -xelatex main.tex
% =====================================================================
\documentclass[aspectratio=169,10pt]{beamer}

% Path has no parent-directory step: Overleaf resolves \input from the project root and
% will not read above it. Locally, the Makefile puts common/ on TEXINPUTS.
% =====================================================================
%  LDA-T318 — Large Language Models and Generative AI for NLP
%  Shared Beamer preamble for all lecture decks.
%  Aarne Talman, University of Helsinki, 2026.
%
%  Engine:  XeLaTeX or LuaLaTeX preferred (gives the Fira font set).
%           pdfLaTeX also works; Metropolis falls back to Latin Modern.
%  Build:   latexmk -xelatex main.tex      (from the deck's own directory)
%
%  Every deck does:
%      \documentclass[aspectratio=169,10pt]{beamer}
%      \input{../common/preamble.tex}
% =====================================================================

% ---------------------------------------------------------------- theme
\usetheme{metropolis}
\metroset{
  titleformat=smallcaps,
  numbering=fraction,
  progressbar=frametitle,
  block=fill,
  sectionpage=progressbar,
  subsectionpage=none,
}

% ---------------------------------------------------------------- packages
% inputenc/fontenc are required under pdfLaTeX and actively harmful under
% XeLaTeX/LuaLaTeX, where Metropolis loads fontspec to get Fira. Guard them.
\usepackage{iftex}
\ifPDFTeX
  \usepackage[utf8]{inputenc}
  \usepackage[T1]{fontenc}
  \usepackage[sfdefault]{FiraSans}
  \usepackage{FiraMono}
\else
  % Metropolis only gets Fira if fontspec can actually find it. TeX Live ships
  % the OTFs but does NOT register them with fontconfig, so a plain
  % \setsansfont{Fira Sans} silently falls back to Latin Modern and the deck
  % stops looking like Metropolis. Load them by filename instead.
  \usepackage{fontspec}
  \setsansfont{FiraSans}[
    Extension      = .otf,
    UprightFont    = *-Book,
    ItalicFont     = *-BookItalic,
    BoldFont       = *-Medium,
    BoldItalicFont = *-MediumItalic]
  \setmonofont{FiraMono}[
    Extension   = .otf,
    UprightFont = *-Regular,
    BoldFont    = *-Medium,
    Scale       = 0.92]
  \renewcommand{\familydefault}{\sfdefault}
\fi
\usepackage[english]{babel}

\usepackage{amsmath,amssymb}
\usepackage{booktabs}
\usepackage{array}
\usepackage{multicol}
\usepackage{xcolor}
\usepackage{graphicx}
\usepackage{tikz}
\usepackage{pgfplots}
\usepackage{listings}
\usepackage{appendixnumberbeamer}

\pgfplotsset{compat=1.17}
\usetikzlibrary{positioning,arrows.meta,fit,backgrounds,calc,matrix,shapes.geometric,decorations.pathreplacing}

% =====================================================================
%  University of Helsinki palette
%  Primary brand colour is UH Hray/black with the "HY blue" accent.
%  Secondary accents chosen for projector contrast, not for print.
% =====================================================================
\definecolor{UHBlue}{HTML}{0C5A9E}   % primary accent
\definecolor{UHDark}{HTML}{15202B}   % body text / dark backgrounds
\definecolor{UHGrey}{HTML}{5A6672}   % secondary text
\definecolor{UHLight}{HTML}{EEF2F6}  % block backgrounds
\definecolor{UHOrange}{HTML}{D1601A} % alerts, gotchas
\definecolor{UHGreen}{HTML}{1F7A5C}  % examples, "this works"
\definecolor{UHRed}{HTML}{A62B2B}    % failure modes
\definecolor{UHSand}{HTML}{C8B88A}   % neutral third accent

\setbeamercolor{normal text}{fg=UHDark, bg=white}
\setbeamercolor{alerted text}{fg=UHOrange}
\setbeamercolor{example text}{fg=UHGreen}
\setbeamercolor{frametitle}{fg=white, bg=UHBlue}
\setbeamercolor{title separator}{fg=UHBlue}
\setbeamercolor{progress bar}{fg=UHBlue}
\setbeamercolor{progress bar in head/foot}{fg=UHBlue}
\setbeamercolor{progress bar in section page}{fg=UHBlue}
\setbeamercolor{structure}{fg=UHBlue}
\setbeamercolor{palette primary}{fg=white, bg=UHBlue}   % section pages, standout frames
\setbeamercolor{block title}{fg=white, bg=UHBlue}
\setbeamercolor{block body}{fg=UHDark, bg=UHLight}
\setbeamercolor{block title alerted}{fg=white, bg=UHOrange}
\setbeamercolor{block body alerted}{fg=UHDark, bg=UHLight}
\setbeamercolor{block title example}{fg=white, bg=UHGreen}
\setbeamercolor{block body example}{fg=UHDark, bg=UHLight}
\setbeamercolor{footnote}{fg=UHGrey}

\setbeamerfont{frametitle}{size=\large, series=\bfseries}
\setbeamerfont{footnote}{size=\tiny}
\setbeamerfont{caption}{size=\scriptsize}

% Frame title bar.
%
% Do NOT reimplement the frametitle template here. Metropolis's own template
% already spans \paperwidth (so `bg=UHBlue' above gives the full-width bar),
% already struts two-line titles without collision, and already draws the
% frametitle progress bar. Reimplementing it means touching internals whose
% names contain `@', which is NOT a letter in a .tex file -- `\@empty' parses
% as `\@' followed by the word "empty" and the build breaks.
%
% To change the bar's height, adjust Metropolis's own padding length instead,
% inside \makeatletter so that `@' is a letter:
%
%   \makeatletter
%   \setlength{\metropolis@frametitle@padding}{2.4ex}  % default 2.2ex
%   \makeatother
%
\makeatletter
\setlength{\metropolis@frametitle@padding}{1.2ex}
\makeatother

% No deck uses \framesubtitle; if one ever does, Metropolis ignores it and the
% subtitle should go in the frame body instead.

% =====================================================================
%  Course metadata — decks override \lecturename and \lecturenumber
% =====================================================================
\newcommand{\coursecode}{LDA-T318}
\newcommand{\coursename}{LLMs and Generative AI for NLP}
\newcommand{\courseteacher}{Aarne Talman}
\newcommand{\courseinstitution}{University of Helsinki}
\providecommand{\lecturename}{}
\providecommand{\lecturedate}{}

% Footline: course code on the left, lecture name centred, n/N on the right.
\setbeamertemplate{frame footer}{\coursecode\quad\textbar\quad\lecturename}

% =====================================================================
%  Semantic macros — use these instead of ad-hoc formatting, so that a
%  change of house style is a single edit here.
% =====================================================================

% --- Sourcing. Every quantitative claim on a slide should carry one. ---
\newcommand{\source}[1]{%
  \vfill{\scriptsize\color{UHGrey}\raggedright Source: #1\par}%
}
\newcommand{\sources}[1]{%
  \vfill{\scriptsize\color{UHGrey}\raggedright Sources: #1\par}%
}
% Inline citation chip, e.g. \cit{Raschka 2026a}
\newcommand{\cit}[1]{{\scriptsize\color{UHGrey}[#1]}}

% --- Counter-intuitive point. Three per lecture; these are exam material. ---
\newenvironment{gotcha}[1][Gotcha]{%
  \begin{alertblock}{#1}%
}{%
  \end{alertblock}%
}

% --- Speculative / unconfirmed claim. Research literacy matters here. ---
\newcommand{\rumour}{{\color{UHOrange}\textbf{[unconfirmed]}}\ }

% --- Key takeaway, for the closing slide of each section. ---
\newenvironment{takeaway}{%
  \begin{block}{Take away}%
}{%
  \end{block}%
}

% --- Live demo marker. ---
\newcommand{\demo}[1]{{\color{UHGreen}\textbf{Demo:}} #1}

% --- Emphasis consistent with the palette. ---
\newcommand{\key}[1]{{\color{UHBlue}\textbf{#1}}}
\newcommand{\warn}[1]{{\color{UHOrange}\textbf{#1}}}
\newcommand{\good}[1]{{\color{UHGreen}\textbf{#1}}}
\newcommand{\bad}[1]{{\color{UHRed}\textbf{#1}}}

% --- Units and quantities used constantly across the decks. ---
\newcommand{\unit}[1]{\,\mathrm{#1}}
\newcommand{\flops}{\mathrm{FLOPs}}
% NB: no \toks macro here --- \toks is a TeX primitive (a token register), so
% \newcommand{\toks} is a hard "already defined" error. No deck used it.
\newcommand{\KV}{\mathrm{KV}}
\newcommand{\softmax}{\operatorname{softmax}}

% =====================================================================
%  Code listings — used for config snippets, Slurm scripts, chat templates.
%  Frames containing these MUST be declared [fragile].
% =====================================================================
\lstdefinestyle{uh}{
  basicstyle=\ttfamily\scriptsize\color{UHDark},
  keywordstyle=\color{UHBlue}\bfseries,
  commentstyle=\color{UHGrey}\itshape,
  stringstyle=\color{UHGreen},
  numberstyle=\tiny\color{UHGrey},
  backgroundcolor=\color{UHLight},
  frame=none,
  xleftmargin=0.6em,
  framexleftmargin=0.6em,
  breaklines=true,
  breakatwhitespace=true,
  showstringspaces=false,
  columns=fullflexible,
  keepspaces=true,
  upquote=true,
  literate={-}{{-}}1 {_}{{\_}}1,
}
\lstset{style=uh}

\lstdefinelanguage{slurm}{
  morekeywords={SBATCH,module,srun,source,export,set,echo},
  morecomment=[l]{\#},
  sensitive=true,
}
\lstdefinelanguage{jinja}{
  morekeywords={for,endfor,if,endif,set,in},
  morecomment=[s]{\{\#}{\#\}},
  sensitive=true,
}

% =====================================================================
%  TikZ styles shared across decks (block diagrams, token grids).
% =====================================================================
\tikzset{
  mod/.style     = {draw=UHBlue, fill=UHLight, rounded corners=2pt,
                    minimum height=7mm, minimum width=22mm,
                    font=\scriptsize, align=center},
  modhi/.style   = {mod, draw=UHOrange, fill=UHOrange!12},
  modmute/.style = {mod, draw=UHGrey, fill=white, text=UHGrey},
  flow/.style    = {-{Latex[length=2mm]}, draw=UHGrey, thick},
  flowhi/.style  = {flow, draw=UHOrange},
  resid/.style   = {-{Latex[length=2mm]}, draw=UHBlue, thick, dashed},
  tokcell/.style = {draw=UHGrey!50, minimum size=4.2mm, font=\tiny, inner sep=0pt},
  tokon/.style   = {tokcell, fill=UHBlue!35},
  tokoff/.style  = {tokcell, fill=UHLight},
  tokmask/.style = {tokcell, fill=UHGrey!15, text=UHGrey!60},
  lbl/.style     = {font=\scriptsize\color{UHGrey}},
}

% =====================================================================
%  Presenter notes.
%
%  Notes are written on every slide. By default they are NOT shown, so
%  `latexmk -xelatex main.tex` produces a clean deck for projection.
%
%  To present with notes on a second screen (pdfpc, Skim, Impressive):
%      uncomment the line below, rebuild to main-notes.pdf
%  To print a notes handout:
%      \setbeameroption{show only notes}
% =====================================================================
% \setbeameroption{show notes on second screen=right}

\setbeamertemplate{note page}{%
  {\scriptsize\color{UHGrey}\coursecode\ \textbar\ \lecturename\par}%
  \vskip0.3em\hrule\vskip0.6em%
  {\small\insertnote}%
}

% ---------------------------------------------------------------- misc
\setbeamertemplate{caption}[numbered]
\setbeamertemplate{itemize items}[circle]
\setlength{\parskip}{0.1em}

% ---------------------------------------------------------------- density
% These decks are text-dense and a 16:9 frame is only 90mm tall, so the
% default vertical rhythm costs several lines per slide. Tightened globally
% here rather than slide by slide.
\setbeamersize{text margin left=0.7cm, text margin right=0.7cm}
\makeatletter
\def\@listi{\leftmargin\leftmargini
  \topsep 1.5\p@ \@plus1\p@ \@minus1\p@
  \parsep 0\p@
  \itemsep 1.5\p@ \@plus1\p@ \@minus1\p@}
\let\@listI\@listi
\makeatother

% The title page is a stack of short paragraphs, so a non-zero \parskip
% pushes it past the bottom of the frame. Neutralise it there only.
\addtobeamertemplate{title page}{\setlength{\parskip}{0pt}}{}

% Tighter lists — the decks are dense.
\setbeamertemplate{itemize/enumerate body begin}{\small}
\setbeamertemplate{itemize/enumerate subbody begin}{\footnotesize}

% Table helpers
\newcolumntype{L}[1]{>{\raggedright\arraybackslash}p{#1}}
\newcolumntype{C}[1]{>{\centering\arraybackslash}p{#1}}
\newcolumntype{R}[1]{>{\raggedleft\arraybackslash}p{#1}}

% Allow slightly overfull boxes without a warning storm
\hfuzz=2pt

% =====================================================================
%  Shared course-map frame.
%  Every lecture opens with this, with the current week highlighted:
%      \coursemapframe{1}   ... \coursemapframe{4}
%  Pass 0 to highlight nothing.
% =====================================================================
% \cmrow{current week}{this row's week}{dates}{Tuesday}{Thursday}{teacher}
% Highlights the row in UH blue bold when the two week numbers match.
% Deliberately avoids \rowcolor so that colortbl is not needed.
% The & and \\ must NOT sit inside the \ifnum branches: TeX scans alignment
% cells while skipping the false branch, so an & there breaks the conditional
% ("Incomplete \ifnum"). Keep the branches down to a macro name and let the
% separators live in the macro bodies, which expand after the test resolves.
\newcommand{\cmhi}[5]{%
  \textcolor{UHBlue}{\textbf{#1}} & \textcolor{UHBlue}{\textbf{#2}} &
  \textcolor{UHBlue}{\textbf{#3}} & \textcolor{UHBlue}{\textbf{#4}} &
  \textcolor{UHBlue}{\textbf{#5}}\\%
}
\newcommand{\cmplain}[5]{#1 & #2 & #3 & #4 & #5\\}
\newcommand{\cmrow}[6]{%
  \ifnum#1=#2\relax\expandafter\cmhi\else\expandafter\cmplain\fi
  {#2}{#3}{#4}{#5}{#6}%
}
\newcommand{\coursemapframe}[1]{%
\begin{frame}{Course overview}
\centering
\small
\setlength{\parskip}{0pt}
\renewcommand{\arraystretch}{1.08}
\begin{tabular}{@{}c L{1.7cm} L{3.1cm} L{3.7cm} L{1.4cm}@{}}
\toprule
\textbf{Wk} & \textbf{Dates} & \textbf{Tuesday} & \textbf{Thursday} & \textbf{Teacher}\\
\midrule
\cmrow{#1}{1}{27 / 29 Oct}{Language models}{Pretraining}{Aarne}
\cmrow{#1}{2}{3 / 5 Nov}{Post-training}{Lab 2: tune vs prompt}{Aarne}
\cmrow{#1}{3}{10 / 12 Nov}{Inference}{Lab 3: serve}{Aarne}
\cmrow{#1}{4}{17 / 19 Nov}{Evaluation}{Lab 4: evaluate}{Jussi}
\cmrow{#1}{5}{24 / 26 Nov}{RAG}{Lab 5: RAG}{Dmitry}
\cmrow{#1}{6}{1 / 3 Dec}{Agents}{Lab 6: agent}{Dmitry}
\cmrow{#1}{7}{8 / 10 Dec}{Course recap}{Special topics (TBC)}{All}
\bottomrule
\end{tabular}

\vspace{1.2em}
{\footnotesize Reusing a task and model across labs is recommended, not required.\\
There are no intermediate submissions; the final report and code repository
determine the grade.}
\note{Fifteen seconds, not more. The point of this slide is orientation, not
content: show them the shape of the course and where today sits inside it.
The final week is deliberately left flexible for a recap and special topics.}
\end{frame}%
}


\renewcommand{\lecturename}{Lecture 1 — Language models}
\renewcommand{\lecturedate}{27 October 2026}

\title{Language models}
\subtitle{Lecture 1: From next-token prediction to modern LLM systems}
\author{Aarne Talman}
\institute{\coursecode\ \textbar\ \coursename\ \textbar\ \courseinstitution}
\date{27 October 2026}

\begin{document}

% =====================================================================
\begin{frame}[plain,shrink=8]
  \titlepage
  \note{Welcome students and briefly introduce the course. By the end, students
  will have trained, tuned, served and evaluated language models on Roihu.

  Keep the administrative introduction brief so there is enough time for the
  conceptual overview.}
\end{frame}

\coursemapframe{1}

\begin{frame}{Today}
  \begin{columns}[T]
    \begin{column}{0.5\textwidth}
      \begin{enumerate}
        \item Course logistics \hfill \textit{10 min}
        \item What is a language model? \hfill \textit{15 min}
        \item From base model to application \hfill \textit{10 min}
        \item Tokenization \hfill \textit{10 min}
      \end{enumerate}
    \end{column}
    \begin{column}{0.5\textwidth}
      \begin{enumerate}
        \setcounter{enumi}{4}
        \item Transformers at a high level \hfill \textit{20 min}
        \item Choosing and reading models \hfill \textit{10 min}
        \item Roihu and Lab 1 \hfill \textit{10 min}
      \end{enumerate}
    \end{column}
  \end{columns}
  \vfill
  \begin{takeaway}
    By the end of today, you should be able to explain how a language model
    turns text into tokens, uses context to predict the next token, and
    generates text one token at a time.
  \end{takeaway}
  \note{Briefly talk through the plan and timings.}
\end{frame}

% =====================================================================
\section{Course logistics}
% =====================================================================

\begin{frame}{The shape of the course}
  \begin{columns}[T]
    \begin{column}{0.56\textwidth}
      Over seven weeks, we cover the main stages in developing and using
      language models:
      \begin{itemize}
        \item architecture and pretraining
        \item post-training
        \item evaluation
        \item inference and serving
        \item retrieval
        \item agents
      \end{itemize}
      \vspace{0.5em}
      \small
      \key{Aarne Talman} --- AMD (weeks 1, 2, 3 and 7)\\
      \key{Jussi Karlgren} --- AMD (week 4)\\
      \key{Dmitry Kan} --- Aiven (weeks 5--6)
    \end{column}
    \begin{column}{0.44\textwidth}
      \begin{block}{Every lab runs on CSC Roihu}
        \small
        NVIDIA GH200 nodes.\\
        One H100-class GPU with 96\,GiB,\\
        plus a 72-core ARM Grace CPU.\\[0.4em]
        Enough memory for the planned labs, including bf16 fine-tuning.
      \end{block}
    \end{column}
  \end{columns}
  \note{Explain how the weeks connect. The adapter from week 2 is served and
  measured in week 3. Week 4 examines how to evaluate models rigorously, and
  the endpoint from week 3 becomes the generator in the week-5 RAG lab.

  Mention the GH200 memory figure now --- 96 GiB --- because it determines
  which fine-tuning and serving configurations fit.}
\end{frame}

\begin{frame}{An optional way to structure your project}
  You may use the same task across several labs. This makes it easier to
  compare prompting, fine-tuning, evaluation, inference and retrieval. You may
  change models when useful.
  \vspace{0.3em}
  \centering
  \renewcommand{\arraystretch}{1.12}
  \begin{tabular}{@{}L{3.0cm} L{2.5cm} L{2.2cm} L{4.3cm}@{}}
    \toprule
    \textbf{Task} & \textbf{Dataset} & \textbf{Metric} & \textbf{Why it works here}\\
    \midrule
    Math problems & GSM8K & exact match & automatically checked answers\\
    \addlinespace[0.35em]
    NLI & XNLI / MNLI & accuracy & multilingual; small label set\\
    \addlinespace[0.35em]
    Finnish QA & TyDi QA (fi) & F1 / EM & connects naturally to RAG\\
    \addlinespace[0.35em]
    Finnish NER & Turku NER & entity F1 & tests structured output\\
    \bottomrule
  \end{tabular}

  \vspace{0.4em}
  \raggedright\small
  You may propose another task if it has \textbf{a labelled test set} and
  \textbf{an automatic metric}. You may also change models between labs when
  that makes sense; explain the reason in your report.
  \note{Emphasise the constraint rather than reading out every option. Each
  lab asks whether a change helped, so students need a metric they can run
  repeatedly.

  The Finnish options connect especially well to the tokenizer and retrieval
  material, but students may choose any task that meets the requirements.}
\end{frame}

\begin{frame}{Choosing models for the labs}
  \begin{columns}[T]
    \begin{column}{0.52\textwidth}
      Suggested models (all open-weight under Apache 2.0):
      \begin{itemize}
        \item Qwen3.5 0.8B, 2B or 4B
        \item Gemma 4 E2B or E4B
      \end{itemize}
      \vspace{0.6em}
      For the week-2 comparison, choose a model that:
      \begin{itemize}
        \item has both \key{base and instruct} versions --- week 2 uses both
        \item fits comfortably on one GH200
        \item has a licence that permits academic use
        \item tokenizes the language of your task reasonably well
      \end{itemize}
    \end{column}
    \begin{column}{0.48\textwidth}
      \begin{block}{If you reuse the same task}
        \footnotesize
        \textbf{Week 1:} choose the task and evaluation set\\
        \textbf{Week 2:} compare prompting with fine-tuning\\
        \textbf{Week 3:} serve and measure a model\\
        \textbf{Week 4:} evaluate the results carefully\\
        \textbf{Weeks 5--6:} add retrieval and tools\\
        \textbf{At the end:} submit the report and code
      \end{block}
      \small
      This is one option, not a requirement.
    \end{column}
  \end{columns}
  \note{Draw attention to the base-and-instruct requirement for week 2.
  Keeping one model across several labs makes comparisons cleaner, but it is
  not required. A change of model can itself be a reasonable methodological
  choice if the student explains it. Larger models may fit on the GPU, but
  they make iteration during a 90-minute lab unnecessarily slow.}
\end{frame}

\begin{frame}{Assessment: one final project}
  The final project accounts for \key{100\% of the course grade}. Submit:
  \begin{itemize}
    \item a \key{final report} with one section for each lecture topic,
          incorporating the related lab work and results
    \item a link to your \key{code repository}, including notebooks, scripts,
          configurations, result files and instructions for running the code
  \end{itemize}

  \vfill
  \begin{block}{No intermediate submissions}
    The weekly lab work contributes material to the final project, but you do
    not submit a separate deliverable after each lab.
  \end{block}

  \vfill
  The report and code are assessed together, with emphasis on sound
  experiments, clear analysis, understanding and reproducibility.
  \vspace{0.3em}

  The final report and repository are due \textbf{31 December}.
  \note{The report should follow the course structure rather than read like
  six disconnected lab logs. Make it explicit that there are no weekly
  submissions: students submit the report and repository at the end.}
\end{frame}

\begin{frame}{Policies}
  \begin{columns}[T]
    \begin{column}{0.5\textwidth}
      \textbf{AI assistants} --- allowed for code and for writing.
      End the report with a short statement of which tools you used and for what.
      \alert{Numbers must come from runs you actually did.}
      \vspace{0.6em}

      \textbf{Pairs} --- may share a project and code. One report, with a short
      statement of who did what.
    \end{column}
    \begin{column}{0.5\textwidth}
      \textbf{Compute etiquette}
      \begin{itemize}
        \item close idle GPU Jupyter sessions
        \item debug on \texttt{gputest} before \texttt{gpumedium}
        \item never run training on a login node
      \end{itemize}
      \vspace{0.4em}
      \textbf{Licences} --- only models and datasets whose licences allow
      academic use. Name them in the report.
    \end{column}
  \end{columns}
  \vfill
  \begin{gotcha}[Please release idle GPUs]
    We are all sharing one course project's compute allocation. A forgotten
    Jupyter session holds a whole GH200 for twelve hours.
  \end{gotcha}
  \note{The important distinction in the AI policy is between assistance with
  code or prose and fabricated experimental evidence. For compute, explain
  that idle interactive sessions consume the shared allocation.}
\end{frame}

% =====================================================================
\section{What is a language model?}
% =====================================================================

\begin{frame}{Start from autocomplete}
  A \key{language model} reads some text and returns a \textbf{probability for
  every possible next word}. For example, after
  \textit{``The capital of Finland is''}:

  \vspace{0.4em}
  \centering
  \begin{tabular}{@{}l r l@{}}
    \toprule
    \textbf{Candidate next word} & \textbf{Probability} & \\
    \midrule
    Helsinki & 0.82 & \good{the model has learned a fact}\\
    a        & 0.06 & plausible English, wrong continuation\\
    Finland  & 0.01 & grammatical, semantically odd\\
    purple   & 0.00001 & \\
    \bottomrule
  \end{tabular}

  \vspace{0.3em}
  \raggedright
  Every number is positive, and across the whole vocabulary they sum to one.
  At each generation step, the model internally produces a distribution like
  this. An API may return only the selected word unless log probabilities are
  requested.

  \vfill
  \begin{takeaway}
    Chat, reasoning, agents and retrieval all build on repeated next-word
    prediction.
  \end{takeaway}
  \note{Ask the room to finish the sentence before showing the table. Use
  their answers to introduce the idea of a distribution over continuations.

  Stress that the output is a \emph{distribution}, not a word. Choosing a word
  is a separate step, covered a few slides from now.

  The numbers are illustrative, not measured. If someone asks, say so.

  Vocabulary note: say ``word'' here and switch to ``token'' once section 3
  earns it. Do not use the two interchangeably before then.}
\end{frame}

\begin{frame}[shrink=12]{From one word to a whole sentence}
  If the model only scores the \emph{next} word, how does it score a whole
  sentence? Multiply, one word at a time.

  \vspace{0.3em}
  \footnotesize
  \[
    P(\text{\textit{The cat sat}})
    = \underbrace{P(\text{The})}_{\text{no context}}
    \times \underbrace{P(\text{cat}\mid\text{The})}_{\text{one word of context}}
    \times \underbrace{P(\text{sat}\mid\text{The cat})}_{\text{two words of context}}
  \]
  \normalsize

  \vspace{0.2em}
  In general, for a sequence of $T$ tokens $x_1,\dots,x_T$:
  \[
    p_\theta(x_1,\dots,x_T) \;=\; \prod_{t=1}^{T} p_\theta\!\left(x_t \mid x_{<t}\right)
  \]
  \emph{In words:} multiply together the probability of each token given
  everything before it. ($x_{<t}$ is everything before position $t$;
  $\theta$ is the model's weights.)

  \vfill
  \begin{block}{Why this matters}
    Splitting a sentence up like this is just the chain rule of probability,
    so no approximation has been introduced. It also turns a sequence of $T$
    tokens into $T$ next-token training examples.
  \end{block}
  \note{Work through the three-word example on the board before showing the
  general form. The chain rule connects sequence probability to the
  next-token prediction task.

  Full chain rule for the sceptics: $P(A,B,C) = P(A)\,P(B\mid A)\,P(C\mid A,B)$,
  which follows from the definition of conditional probability applied twice.
  It holds for any ordering of the variables; left-to-right is simply the
  ordering that matches how text is written and read.

  If useful, note that masked language models and diffusion language models
  use different objectives; ``language model'' is broader than
  ``autoregressive decoder''.}
\end{frame}

\begin{frame}[shrink=12]{The training objective}
  Training adjusts the weights so that \textbf{real text gets high
  probability}. Written as something to minimise:
  \[
    \mathcal{L} \;=\; -\frac{1}{T}\sum_{t=1}^{T}\log p_\theta\!\left(x_t \mid x_{<t}\right)
  \]
  \emph{In words:} at every position, look at the probability the model gave
  the word that actually came next, and penalise it for being small.

  \vspace{0.5em}
  \begin{columns}[T]
    \begin{column}{0.52\textwidth}
      \begin{block}{Why the logarithm}
        \small
        It turns the product on the previous slide into a sum, which is
        numerically stable and gives a per-token average that does not depend
        on sentence length.
      \end{block}
    \end{column}
    \begin{column}{0.48\textwidth}
      \begin{block}{What it is used for}
        \small
        \key{Cross-entropy loss} provides the training signal used to update
        the model's weights. It is the same loss used in classification, with
        the vocabulary entries acting as classes.
      \end{block}
    \end{column}
  \end{columns}

  \vfill
  \begin{takeaway}
    The same basic objective is used to train small classroom models and large
    production models.
  \end{takeaway}
  \note{Connect explicitly to what they already know: this is softmax
  classification over a 50k-way vocabulary, run once per position. Most of
  them have trained exactly this loss on MNIST or a sentiment classifier.

  If a perfect model assigned probability 1 to every correct token, the loss
  would be 0. Random guessing over a 50k vocabulary gives
  $\log 50000 \approx 10.8$ nats. Real models land around 2--3. Those three
  numbers are the whole scale and they set up perplexity two slides on.

  Units: nats if $\log$ is natural, bits if base 2. Papers use both; divide by
  $\ln 2$ to convert.}
\end{frame}

\begin{frame}[shrink=8]{From counts to learned representations}
  One way to estimate these probabilities is to \textbf{count}: take a large
  corpus and count how often ``sat'' follows ``The cat''. This gives an
  \key{n-gram model}. N-grams remained strong baselines into the early 2010s
  as neural language models increasingly replaced them.

  \vspace{0.4em}
  \begin{columns}[T]
    \begin{column}{0.5\textwidth}
      \begin{gotcha}[Where counting struggles]
        Most contexts you meet at test time were \textbf{never seen in
        training}, so the count is zero and the probability is zero.

        \vspace{0.3em}
        And counts cannot share what they learn. ``The cat sat'' tells the
        model nothing about ``The dog sat'', because the two contexts are
        just different strings.
      \end{gotcha}
    \end{column}
    \begin{column}{0.5\textwidth}
      \begin{block}{What vectors fix}
        \small
        Each word becomes a \key{vector}, and similar words get similar
        vectors, so ``dog'' inherits statistics from ``cat'' for free.

        \vspace{0.3em}
        This sharing lets the model assign sensible scores to contexts that
        never appeared verbatim in the training data. This ability to
        generalise beyond exact matches is why learned representations matter.
      \end{block}
    \end{column}
  \end{columns}

  \vfill
  \raggedright
  \footnotesize
  \textbf{The lineage:} n-gram $\to$ neural LM (Bengio 2003) $\to$ RNN/LSTM
  $\to$ \textbf{attention}, invented to fix the RNN bottleneck in translation
  $\to$ \textbf{transformer}, which kept attention and removed the recurrence.
  \source{Bengio et al., \textit{A Neural Probabilistic Language Model}, 2003; Vaswani et al., \textit{Attention Is All You Need}, 2017.}
  \note{Point out that attention predates the transformer. The transformer
  retained attention while removing recurrence, which made sequence training
  much easier to parallelise.

  Numbers if they want them: a 50k vocabulary has $50000^3 \approx 10^{14}$
  possible trigrams, against maybe $10^{10}$ words of training text. The
  overwhelming majority of contexts are unseen, which is why the field spent
  two decades on smoothing methods (Katz back-off, Kneser--Ney) that were
  entirely about redistributing probability mass to unseen events.

  Embeddings provide a different kind of generalisation: similar contexts can
  share statistical strength.}
\end{frame}

\begin{frame}{From probabilities to text: where ``generative'' comes in}
  A trained model gives you a distribution. To get \emph{text} you have to
  \textbf{choose}, and then feed that choice back in as context.

  \vspace{0.4em}
  \centering
  \begin{tikzpicture}[node distance=11mm, font=\scriptsize,
      box/.style={draw=UHBlue, rounded corners, inner sep=4pt,
                  minimum height=8mm, align=center}]
    \node[box] (ctx)  {context so far};
    \node[box, right=of ctx] (model) {language model};
    \node[box, right=of model] (logits) {logits};
    \node[box, right=of logits] (probs) {probabilities};
    \node[box, right=of probs] (pick) {choose a token};
    \draw[->, UHBlue, thick] (ctx) -- (model);
    \draw[->, UHBlue, thick] (model) -- (logits);
    \draw[->, UHBlue, thick] (logits) -- node[above]{softmax} (probs);
    \draw[->, UHBlue, thick] (probs) -- (pick);
    \draw[->, UHOrange, thick] (pick.south) -- ++(0,-5mm) -| node[below, pos=0.25]{append, repeat} (ctx.south);
  \end{tikzpicture}

  \vspace{0.5em}
  \raggedright
  \begin{columns}[T]
    \begin{column}{0.48\textwidth}
      \textbf{Always take the most likely word} --- repeatable, and prone to
      loops and flat prose.
    \end{column}
    \begin{column}{0.52\textwidth}
      \textbf{Sample from the distribution} --- varied and more natural;
      \key{temperature} flattens or sharpens it first. Same prompt, different
      answer each time.
    \end{column}
  \end{columns}

  \vfill
  \begin{takeaway}
    If the same prompt gives two different answers, sampling is usually the
    reason. In week 3, we look at how systems make this loop fast.
  \end{takeaway}
  \note{The raw scores are called \emph{logits}; softmax exponentiates and
  normalises them into a distribution. Temperature $T$ divides the logits
  before the softmax: $T<1$ sharpens towards greedy, $T>1$ flattens towards
  uniform, $T\to 0$ is greedy. Top-$k$ and top-$p$ truncate the tail first.
  Sampling methods are covered in week 3; here, the distinction between
  scoring and choosing is enough.

  Two consequences worth stating out loud now, because both will bite them in
  the labs:
  (1) the same prompt gives different answers, so any single generation is
  weak evidence --- this is Jussi's week-3 problem;
  (2) generation is sequential by construction, one forward pass per token,
  which is the entire reason inference is expensive. Flag that this is the
  thing you personally work on, and that even at temperature 0 you do not
  quite get determinism --- week 3.}
\end{frame}

\begin{frame}[shrink=8]{Perplexity}
  Cross-entropy is useful for training, but its scale is hard to interpret.
  \key{Perplexity} turns it into an approximate ``number of plausible choices'':
  \[
    \mathrm{PPL} \;=\; e^{\mathcal{L}}
  \]
  For example, if the average loss is $\mathcal{L}=2$, then
  $\mathrm{PPL}=e^2\approx 7.4$. The model is about as uncertain as a uniform
  choice among seven or eight tokens at each step.
  \vspace{0.3em}

  \textbf{Lower is better.} Perfect prediction gives $\mathrm{PPL}=1$.
  \vspace{0.3em}

  \textbf{What it is used for:} tracking progress during training and
  comparing checkpoints on the same tokenized validation data.
  \vfill
  \begin{gotcha}[Compare perplexity only with the same tokenizer]
    Changing the tokenizer changes what counts as one prediction. A larger
    vocabulary produces fewer but harder tokens, so its perplexity is on a
    different scale. Use \emph{bits per byte} for comparisons across
    tokenizers.
  \end{gotcha}
  \note{Start with the example rather than the formula. Perplexity does not
  mean the model literally considers exactly 7.4 tokens; it describes an
  equivalent level of average uncertainty.

  The tokenizer warning matters for Lab 1: different vocabulary sizes change
  what counts as one prediction, so their perplexities are not directly
  comparable. This is why the lab asks for a loss-versus-compute plot.}
\end{frame}

\begin{frame}{What next-token prediction can—and cannot—teach}
  \begin{columns}[T]
    \begin{column}{0.52\textwidth}
      \begin{block}{Why useful abilities emerge}
        To predict text well, a model must learn recurring patterns in
        language: grammar, meaning, facts, discourse, code and common forms of
        reasoning.
        \vspace{0.4em}

        The better it represents these patterns, the better it can predict
        what comes next.
      \end{block}
      \begin{gotcha}[What the objective does not guarantee]
        Frequent text is not necessarily true or useful. A fluent continuation
        is therefore not evidence that a claim is correct.
      \end{gotcha}
    \end{column}
    \begin{column}{0.48\textwidth}
      \begin{block}{A simple example}
        \footnotesize
        \textit{``The Helsinki train was delayed, so Anna arrived''}
        \textbf{\textcolor{UHBlue}{late}}\\[0.4em]
        Predicting ``late'' requires more than the previous word. The model
        needs to connect the delay, the train and Anna across the sentence.
        \vspace{0.4em}

        Repeating this task over a large corpus encourages the model to learn
        such relationships.
      \end{block}
      \small
      A base model learns to \emph{continue text}; it is not automatically
      trained to follow a user's instructions. That is the topic of week 2.
    \end{column}
  \end{columns}
  \note{Read the example aloud and pause before ``late''. Ask what information
  students used to complete it. The model encounters similar prediction
  problems throughout the corpus, which creates pressure to represent
  linguistic and factual patterns.

  Keep the limitation equally clear: likelihood and truth are different
  objectives. Week 2 explains how post-training changes behaviour.}
\end{frame}

\begin{frame}{What comes with a model release?}
  \begin{columns}[T]
    \begin{column}{0.48\textwidth}
      \begin{itemize}
        \item \key{weights} --- the \texttt{.safetensors} files
        \item \key{tokenizer} --- vocabulary plus merge rules
        \item \key{chat template} --- a Jinja string that turns a conversation
              into one flat sequence
        \item \key{generation config} --- default sampling parameters
        \item \key{inference code} --- the architecture itself
      \end{itemize}
    \end{column}
    \begin{column}{0.52\textwidth}
      \begin{gotcha}
        Three of those five are \textbf{not} in the weights file, and each of
        them can silently ruin your evaluation.
        \vspace{0.3em}

        The most common failure in Lab 2: applying a chat template to a
        \textbf{base} model, which has never seen those tokens.
      \end{gotcha}
    \end{column}
  \end{columns}
  \vspace{0.3em}
  \small
  Browse open-weight model releases at
  \href{https://huggingface.co/models}{\texttt{huggingface.co/models}}.
  \vfill
  \begin{takeaway}
    Model behaviour depends on the weights and on how text is tokenized,
    formatted and decoded.
  \end{takeaway}
  \note{Return to this list in weeks 2 and 4. Unexpected output often comes
  from a tokenizer, chat-template or generation-setting mismatch rather than
  from the weights themselves.}
\end{frame}

\begin{frame}{The life cycle of a language model}
  \centering
  \begin{tikzpicture}[node distance=5mm and 4mm]
    \node[mod] (data) {text and code};
    \node[mod, right=of data] (tok) {tokenizer};
    \node[modhi, right=of tok] (pre) {pretraining};
    \node[mod, right=of pre] (base) {base model};
    \node[modhi, below=of base] (post) {post-training};
    \node[mod, left=of post] (inst) {instruct /\\reasoning model};
    \node[mod, left=of inst] (app) {application};
    \draw[flow] (data) -- (tok);
    \draw[flow] (tok) -- (pre);
    \draw[flow] (pre) -- (base);
    \draw[flow] (base) -- (post);
    \draw[flow] (post) -- (inst);
    \draw[flow] (inst) -- (app);
  \end{tikzpicture}

  \vspace{0.8em}
  \begin{columns}[T]
    \begin{column}{0.33\textwidth}
      \textbf{Pretraining}\\
      Learn broad patterns from a large corpus by predicting the next token.
    \end{column}
    \begin{column}{0.33\textwidth}
      \textbf{Post-training}\\
      Teach the base model to follow instructions, use tools or reason in a
      desired format.
    \end{column}
    \begin{column}{0.33\textwidth}
      \textbf{Application}\\
      Add prompting, retrieval, tools, evaluation and an inference system.
    \end{column}
  \end{columns}
  \note{Use this as the map for the whole course. The underlying next-token
  objective remains central, but the data and training signal change. An LLM
  application also includes substantial code and data outside the model.}
\end{frame}

\begin{frame}[shrink=2]{The same model can play different roles}
  \begin{columns}[T]
    \begin{column}{0.48\textwidth}
      \begin{block}{Base model}
        Continues text in the style of its training data.
        \vspace{0.3em}

        Useful for studying pretraining and as a starting point for further
        training, but not designed as a chat assistant.
      \end{block}
      \begin{block}{Instruct model}
        Post-trained on instructions and preferences. Expects a particular
        chat template and usually follows user requests more directly.
      \end{block}
    \end{column}
    \begin{column}{0.48\textwidth}
      \begin{block}{Reasoning model}
        Trained to spend additional tokens on difficult problems before
        producing an answer.
      \end{block}
      \begin{block}{Model inside a system}
        Retrieval can supply documents; tools can supply calculations or
        external actions. The model remains the text-prediction component.
      \end{block}
    \end{column}
  \end{columns}
  \vfill
  Week 2 covers post-training; weeks 5 and 6 add retrieval and tools.
  \note{Students often use ``the model'' to refer to an entire product. Keep
  the distinction between model behaviour learned in the weights and
  capabilities supplied by the surrounding system.}
\end{frame}

% =====================================================================
\section{Tokenization}
% =====================================================================

\begin{frame}{Characters, words or subwords?}
  \centering
  \begin{tabular}{@{}l L{4.2cm} L{4.6cm}@{}}
    \toprule
    & \textbf{Advantage} & \textbf{Trade-off}\\
    \midrule
    \textbf{characters} & tiny vocabulary, no OOV & sequences 4--5$\times$ longer;
      attention is $O(n^2)$ \\
    \textbf{words} & short sequences, interpretable & open vocabulary; every typo,
      compound and inflection is OOV \\
    \textbf{subwords} & bounded vocabulary, no OOV, reasonable length & the
      segmentation depends on the training data \\
    \bottomrule
  \end{tabular}
  \vspace{1.0em}

  Byte-level BPE takes the subword route and removes the last OOV case:
  the base alphabet is the \key{256 byte values}, so text in any language,
  including emoji, can be represented through its bytes.
  \note{Connect sequence length to the attention cost discussed later.
  Character models are possible, but their longer sequences make standard
  attention more expensive.

  Byte-level tokenizers generally avoid \texttt{<unk>} for text.}
\end{frame}

\begin{frame}{Byte-level BPE, worked}
  Start with individual bytes. Find a frequent adjacent pair, merge it into a
  new token, and record the merge. Repeat until the vocabulary is large enough.
  \vspace{0.3em}
  \centering
  \begin{tabular}{@{}c L{4.0cm} L{4.8cm}@{}}
    \toprule
    \textbf{Step} & \textbf{Merge learned} & \textbf{A word after the merge}\\
    \midrule
    0 & none yet & \texttt{l o w e s t}\\
    \addlinespace[0.3em]
    1 & \texttt{l} + \texttt{o} $\rightarrow$ \texttt{lo}
      & \texttt{lo w e s t}\\
    \addlinespace[0.3em]
    2 & \texttt{lo} + \texttt{w} $\rightarrow$ \texttt{low}
      & \texttt{low e s t}\\
    \addlinespace[0.3em]
    3 & \texttt{e} + \texttt{s} $\rightarrow$ \texttt{es}
      & \texttt{low es t}\\
    \addlinespace[0.3em]
    4 & \texttt{es} + \texttt{t} $\rightarrow$ \texttt{est}
      & \texttt{low est}\\
    \bottomrule
  \end{tabular}

  \vspace{0.3em}
  \raggedright
  This is one possible merge sequence learned from a tiny corpus containing
  words such as \texttt{low}, \texttt{lower}, \texttt{newest} and
  \texttt{widest}. A real tokenizer learns thousands of merges.
  \vspace{0.1em}

  \textbf{Why it matters:} these merges determine the units the model learns
  to represent and how many tokens each input consumes.
  \note{Walk through the rightmost column from top to bottom. Explain that the
  training corpus determines which pair is selected at each step. Encoding
  later replays the learned merges in order.}
\end{frame}

\begin{frame}{Choosing the vocabulary size}
  \begin{columns}[T]
    \begin{column}{0.48\textwidth}
      \begin{block}{What a larger vocabulary buys}
      \begin{itemize}
        \item common words and word parts fit into fewer tokens
        \item sequences become shorter
        \item more text fits into the same context window
        \item attention has fewer positions to process
      \end{itemize}
      \end{block}
    \end{column}
    \begin{column}{0.48\textwidth}
      \begin{gotcha}[What it costs]
        Every vocabulary item needs an input embedding and an output score.
        A larger vocabulary therefore uses more model parameters.
        \vspace{0.4em}

        Rare tokens also receive fewer training examples than reusable
        subword pieces.
      \end{gotcha}
    \end{column}
  \end{columns}
  \vfill
  \begin{block}{Decoder-only examples—not target values}
    \centering\small
    \begin{tabular}{@{}C{2.8cm} C{2.8cm} C{2.8cm} C{2.8cm}@{}}
      GPT-2 (English) & Poro (Finnish + English) &
      Qwen2.5 (multilingual) & Gemma 2 (multilingual)\\
      \textbf{50.3k} & \textbf{128k} & \textbf{151.9k} & \textbf{256k}
    \end{tabular}
  \end{block}
  \source{Model tokenizer/config files; Luukkonen et al.\ 2024 (Poro).}
  \note{In Lab 1 they train at two or three vocabulary sizes and will see this
  trade directly: the bigger vocabulary gives shorter sequences but spends
  more of a tiny parameter budget on embeddings. These examples illustrate
  scale; vocabulary size alone does not determine how well a language is
  represented.}
\end{frame}

\begin{frame}{Fertility: why some languages use more tokens}
  \key{Fertility} = tokens per word. The same sentence, the same tokenizer,
  different languages:
  \vspace{0.5em}
  \centering
  \begin{tabular}{@{}l c L{7.0cm}@{}}
    \toprule
    \textbf{Language} & \textbf{Typical fertility} & \textbf{Why}\\
    \midrule
    English & $\approx$ 1.2--1.3 & dominates the training corpus; common words are single tokens\\
    German & $\approx$ 1.6--2.0 & compounding\\
    Finnish & $\approx$ 2.0--2.8 & many endings create many word forms; less Finnish in the training data\\
    Thai, Georgian, Amharic & 3--5+ & little training data, non-Latin script\\
    \bottomrule
  \end{tabular}

  \vspace{0.9em}
  \raggedright
  \textit{talossanikinko} $\to$ a single Finnish word; several tokens in every
  mainstream tokenizer.
  \vspace{0.3em}

  \textbf{Lab 1, Part A:} measure this yourself on FLORES-200, for your own
  tokenizer and for Qwen3.5's and Gemma 4's.
  \note{Put a real Finnish sentence through two tokenizers live if the
  network cooperates; the visual of the same sentence becoming twice as many
  coloured boxes is worth more than the table.

  Exact fertility numbers depend on the tokenizer and the corpus --- present
  these as ranges, and tell them their lab numbers are the real ones.}
\end{frame}

\begin{frame}{What high fertility costs you}
  \begin{columns}[T]
    \begin{column}{0.33\textwidth}
      \begin{block}{Cost}
        APIs bill per token. The same content may cost more to process in
        Finnish than in English.
      \end{block}
    \end{column}
    \begin{column}{0.33\textwidth}
      \begin{block}{Context}
        A fixed token window can hold substantially less Finnish text than
        English text.
      \end{block}
    \end{column}
    \begin{column}{0.33\textwidth}
      \begin{block}{Quality}
        More tokens per unit of meaning can make dependencies longer and
        spread the available training signal across more tokens.
      \end{block}
    \end{column}
  \end{columns}
  \vfill
  \begin{gotcha}[The tokenizer is usually fixed before pretraining]
    The tokenizer is trained first and normally kept fixed for the rest of
    training. Fine-tuning cannot fully compensate for poor segmentation.
  \end{gotcha}
  \note{Connect tokenization to language coverage. Morphology, script and the
  language distribution in the tokenizer's training corpus all affect
  fertility.

  Ask students which languages they expect to use and whether the same issue
  might affect their projects.}
\end{frame}

% =====================================================================
\section{How transformers process context}
% =====================================================================

\begin{frame}{A transformer language model, end to end}
  \centering
  \begin{tikzpicture}[node distance=8mm and 12mm,
      box/.style={draw=UHBlue, rounded corners, align=center,
                  inner sep=5pt, minimum height=10mm, minimum width=25mm}]
    \node[box] (text) {text};
    \node[box, right=of text] (tokens) {token\\IDs};
    \node[box, right=of tokens] (embed) {token\\vectors};
    \node[box, below=of tokens, draw=UHOrange] (blocks) {transformer\\blocks};
    \node[box, right=of blocks] (logits) {vocabulary\\scores};
    \node[box, right=of logits] (next) {next\\token};
    \draw[->, UHBlue, thick] (text) -- node[above, font=\scriptsize]{tokenize} (tokens);
    \draw[->, UHBlue, thick] (tokens) -- node[above, font=\scriptsize]{embed} (embed);
    \draw[->, UHBlue, thick] (embed.south) -- ++(0,-4mm) -| (blocks.north);
    \draw[->, UHBlue, thick] (blocks) -- (logits);
    \draw[->, UHBlue, thick] (logits) -- node[above, font=\scriptsize]{choose} (next);
  \end{tikzpicture}

  \vspace{0.5em}
  Each transformer block alternates between:
  \begin{columns}[T]
    \begin{column}{0.48\textwidth}
      \begin{block}{Attention}
        Lets each token gather relevant information from other positions.
      \end{block}
    \end{column}
    \begin{column}{0.48\textwidth}
      \begin{block}{Feed-forward network}
        Transforms the information at each position independently.
      \end{block}
    \end{column}
  \end{columns}
  \vfill
  Stack the block many times, then map the final vector to one score for every
  vocabulary item.
  \note{Give students the full data flow before introducing queries, keys and
  values. Attention is only one part of the block; the repeated stack produces
  the contextual representation used for next-token prediction.}
\end{frame}

\begin{frame}{Why transformers use attention}
  Language depends on context. In \emph{``The animal did not cross the road
  because it was tired''}, understanding \emph{it} requires information from
  several words earlier.
  \vspace{0.5em}
  \begin{columns}[T]
    \begin{column}{0.48\textwidth}
      \begin{block}{Earlier recurrent models}
        An RNN carries context forward through one hidden state, one token at a
        time.
        \vspace{0.3em}

        Information travels through many sequential steps, and the positions
        cannot be computed in parallel during training.
      \end{block}
    \end{column}
    \begin{column}{0.48\textwidth}
      \begin{block}{Self-attention}
        Each position can retrieve information directly from relevant earlier
        positions.
        \vspace{0.3em}

        During training, all positions can be computed in parallel; a causal
        mask prevents access to future tokens.
      \end{block}
    \end{column}
  \end{columns}
  \vfill
  Inside a transformer block, \key{attention} exchanges information between
  positions; the \key{feed-forward network} transforms each position
  independently.
  \note{Use the pronoun example to motivate the need for information from
  another position. Then contrast recurrence with direct access and parallel
  training. The next slides explain how attention performs the retrieval and
  how the causal mask makes parallel training valid.}
\end{frame}

\begin{frame}{Queries, keys and values}
  Think of attention as a learned search over the context. Each token produces:
  \vspace{0.3em}
  \begin{columns}[T]
    \begin{column}{0.31\textwidth}
      \begin{block}{Query}
        \centering
        What information am I looking for?
      \end{block}
    \end{column}
    \begin{column}{0.31\textwidth}
      \begin{block}{Key}
        \centering
        What information do I contain?
      \end{block}
    \end{column}
    \begin{column}{0.31\textwidth}
      \begin{block}{Value}
        \centering
        What information should I pass on?
      \end{block}
    \end{column}
  \end{columns}
  \vspace{0.2em}

  \centering
  \begin{tikzpicture}[node distance=12mm,
      box/.style={draw=UHBlue, rounded corners, align=center,
                  minimum width=38mm, minimum height=12mm, inner sep=4pt}]
    \node[box] (match) {compare query\\with every key};
    \node[box, right=of match] (weights) {softmax\\match weights};
    \node[box, right=of weights] (blend) {weighted blend\\of values};
    \draw[flow] (match) -- (weights);
    \draw[flow] (weights) -- (blend);
  \end{tikzpicture}

  \vspace{0.1em}
  \[
    \mathrm{Attention}(Q,K,V)
    = \softmax\!\left(\frac{QK^{\top}}{\sqrt{d_k}}\right)V
  \]

  \raggedright
  \textbf{What it is used for:} letting each token retrieve the parts of the
  context that are most useful for its current representation.
  \note{In the library analogy, the query is the request, the keys describe
  what is on the shelves, and the values are the contents retrieved. Softmax
  produces a weighted blend rather than choosing exactly one item.

  Projections, if you want them on the board: for each token's vector $x_t$,
  $q_t = W_Q x_t$, $k_t = W_K x_t$, $v_t = W_V x_t$, with $W_Q, W_K, W_V$
  learned. Stacking over $t$ gives $Q$, $K$, $V$.

  The $\sqrt{d_k}$ argument is one line and worth deriving if the room is
  following: for $q, k$ with independent unit-variance entries, the dot
  product $q\cdot k$ is a sum of $d_k$ such terms, so it has variance $d_k$
  and standard deviation $\sqrt{d_k}$. Dividing restores unit scale. If you
  skip the division, the softmax inputs grow with model width, the softmax
  becomes nearly one-hot, and the gradients vanish.

  The causal mask is deliberately not on this slide --- it is the next one.
  If someone asks where it went, say ``next slide'' and move on.}
\end{frame}

\begin{frame}{The causal mask}
  \begin{columns}[T]
    \begin{column}{0.42\textwidth}
      \centering
      \begin{tikzpicture}[scale=1.1]
        \foreach \r in {1,...,7}{
          \foreach \c in {1,...,7}{
            \ifnum\c>\r\relax
              \node[tokmask] at (\c*0.52, -\r*0.52) {};
            \else
              \node[tokon] at (\c*0.52, -\r*0.52) {};
            \fi
          }
        }
        \node[lbl, anchor=south] at (2.1,-0.25) {keys $\rightarrow$};
        \node[lbl, rotate=90, anchor=south] at (0.1,-2.1) {$\leftarrow$ queries};
      \end{tikzpicture}
    \end{column}
    \begin{column}{0.58\textwidth}
      Each row shows which earlier tokens one position is allowed to use.
      \vspace{0.4em}

      The grey upper triangle is the \key{causal mask}: future positions are
      hidden.
      \vspace{0.4em}

      \textbf{Why it is used:} during training, we can compute all positions
      \textbf{in parallel} without letting a token see the answer it should
      predict.
      \vspace{0.4em}

      During generation, each new token adds one row. This sequential step is
      one reason generation is slower than training.
    \end{column}
  \end{columns}
  \note{Connect the mask to parallel training. The whole sequence provides
  training targets in one forward pass while the mask preserves the
  left-to-right prediction task.

  The attention matrix is an intermediate result, not a learned parameter.
  FlashAttention avoids storing the full matrix in GPU memory.

  Forward-reference week 3: prefill does the whole triangle at once and decode
  does one row at a time, and most of the inference lecture follows from
  that.}
\end{frame}

\begin{frame}{Multiple heads}
  Multiple heads let the model retrieve several kinds of contextual
  information at once. Run $h$ attentions in parallel with \emph{different}
  learned projections, each of width
  $d_{\text{head}} = d_{\text{model}}/h$, then concatenate and project.
  \vspace{0.5em}
  \begin{columns}[T]
    \begin{column}{0.52\textwidth}
      \begin{itemize}
        \item Same cost as one wide head, but $h$ separate ways to route information
        \item Empirically heads specialise: local dependencies, syntactic
          relations, long-range semantic links, positional patterns
        \item Many heads turn out to be prunable after training --- a hint that
          this is redundancy as much as specialisation
      \end{itemize}
    \end{column}
    \begin{column}{0.48\textwidth}
      \begin{block}{The inference-time cost}
        Every head keeps its own $K$ and $V$ for every token in the context,
        for every layer, for every request in the batch.
        \vspace{0.4em}
        That stored state is the \key{KV cache}.
      \end{block}
    \end{column}
  \end{columns}
  \note{Avoid treating each head as a cleanly interpretable feature. Some
  heads show recognisable patterns, but heads can also be redundant and their
  roles need not have a simple linguistic description.

  Use the KV-cache box to introduce the next slide's memory calculation.}
\end{frame}

\begin{frame}{Why attention becomes expensive}
  \begin{columns}[T]
    \begin{column}{0.48\textwidth}
      \begin{block}{Compute: compare token pairs}
        For $n$ tokens, $QK^{\top}$ contains $n^2$ scores.
        \vspace{0.4em}

        Double the context length and attention performs roughly
        \textbf{four times as much} pairwise work.
      \end{block}
    \end{column}
    \begin{column}{0.48\textwidth}
      \begin{block}{Memory: keep the past}
        During generation, the model stores earlier keys and values so it
        does not recompute them at every step.
        \vspace{0.4em}

        This \key{KV cache} grows with the context length and the number of
        simultaneous requests.
      \end{block}
    \end{column}
  \end{columns}
  \vfill
  \begin{takeaway}
    Compute limits how much context one request can process. KV-cache memory
    often limits how many requests a server can handle at once.
  \end{takeaway}
  \note{Keep the two costs separate. Quadratic compute matters especially
  during training and prompt processing. KV-cache memory matters especially
  during autoregressive generation and concurrent serving.}
\end{frame}

\begin{frame}{KV cache: a worked example}
  For a 4B-class model with 36 layers, 8 KV heads, head width 128 and bf16
  values:
  \[
    \underbrace{36}_{\text{layers}} \times
    \underbrace{8 \times 128}_{\text{width}} \times
    \underbrace{2}_{\text{keys \emph{and} values}} \times
    \underbrace{2}_{\text{bytes}}
    \;\approx\; 147\unit{kB}\ \textbf{per token}
  \]

  \vspace{0.5em}
  \centering
  \begin{tikzpicture}[node distance=10mm,
      box/.style={draw=UHBlue, rounded corners, align=center,
                  minimum width=35mm, minimum height=14mm, inner sep=4pt}]
    \node[box] (token) {\textbf{147\,kB}\\per token};
    \node[box, right=of token] (request) {\textbf{1.2\,GB}\\one 8k-token request};
    \node[box, right=of request] (batch) {\textbf{39\,GB}\\32 such requests};
    \draw[flow] (token) -- node[above, font=\scriptsize]{$\times 8{,}000$} (request);
    \draw[flow] (request) -- node[above, font=\scriptsize]{$\times 32$} (batch);
  \end{tikzpicture}

  \vfill
  \raggedright
  The model's bf16 weights take about \textbf{8\,GB}; in this example, the
  cache takes much more. This is why serving systems spend so much effort on
  KV-cache management.
  \note{Work through the arithmetic with the class. The important comparison
  is between the model weights and the cache needed for concurrent users.

  In this example, the \emph{weights} of a 4B model take about 8 GB in bf16,
  while the cache for 32 users takes about 39 GB. The KV cache can therefore
  determine how many requests fit on a GPU.

  General form, if you want it on the board --- KV bytes equals
  batch $\times$ tokens $\times$ layers $\times$ KV heads $\times$ head
  dimension $\times$ 2 (for K and V) $\times$ bytes per value. The slide fixes
  batch and tokens to 1 to get the per-token figure, then scales up.

  Why ``two bytes'': bf16. Week 3 covers what happens when you drop to 8 or
  even 4 bits for the cache.}
\end{frame}

% Detailed architecture material retained as reference, but omitted from the
% main Lecture 1 path. Lecture 1 is an overview; these mechanisms can be used
% in a later architecture or inference lecture.
\iffalse

% =====================================================================
\section{The 2026 decoder block}
% =====================================================================

\begin{frame}[shrink=10]{Anatomy of one block}
  \begin{columns}[T]
    \begin{column}{0.42\textwidth}
      \centering
      \begin{tikzpicture}[node distance=2.4mm]
        \node[modmute] (in) {input $x$};
        \node[mod, below=of in] (n1) {RMSNorm};
        \node[modhi, below=of n1] (at) {Attention\\{\tiny GQA / MLA / SWA}};
        \node[mod, below=of at] (a1) {$+$ residual};
        \node[mod, below=of a1] (n2) {RMSNorm};
        \node[modhi, below=of n2] (ff) {FFN\\{\tiny SwiGLU, dense or MoE}};
        \node[mod, below=of ff] (a2) {$+$ residual};
        \node[modmute, below=of a2] (out) {output};
        \draw[flow] (in) -- (n1);
        \draw[flow] (n1) -- (at);
        \draw[flow] (at) -- (a1);
        \draw[flow] (a1) -- (n2);
        \draw[flow] (n2) -- (ff);
        \draw[flow] (ff) -- (a2);
        \draw[flow] (a2) -- (out);
        \draw[resid] (in.west) -- ++(-0.75,0) |- (a1.west);
        \draw[resid] (a1.east) -- ++(0.75,0) |- (a2.east);
      \end{tikzpicture}
    \end{column}
    \begin{column}{0.58\textwidth}
      \textbf{Pre-norm.} Normalise \emph{before} each sublayer rather than
      after. This gives gradients a cleaner path through the residual stream
      and generally makes deep models easier to train.
      \vspace{0.4em}

      \textbf{Two sublayers.} Attention mixes positions, the FFN mixes
      channels, and each sits inside a residual connection.
      \vspace{0.4em}

      \textbf{Repeat this block} many times, then apply a final norm and the
      output projection back to vocabulary size.
      \vspace{0.4em}

      Many recent architecture changes modify one of the two orange boxes.
    \end{column}
  \end{columns}
  \note{Draw attention to the dashed residual lines. The residual stream is
  the spine of the model; the sublayers read from it and write back into it.
  That mental model pays off when we get to mHC later, which modifies the
  residual path itself.

  ``Everything new is in the orange boxes'' is the organising claim for the
  next ten minutes.}
\end{frame}

\begin{frame}{Normalisation}
  \begin{columns}[T]
    \begin{column}{0.5\textwidth}
      \key{Normalisation} keeps activations at a manageable scale and helps
      stabilise training.
      \vspace{0.4em}

      \textbf{LayerNorm $\to$ RMSNorm.} Divide each vector by its own
      root-mean-square size, then rescale by a learned gain. Compared with the
      older LayerNorm it drops the mean subtraction and the bias. It is
      cheaper and is now common in decoder-only language models.
      \vspace{0.5em}

      \textbf{QK-Norm} --- do the same to the queries and keys \emph{inside}
      attention. This can reduce loss spikes during training and is used in
      many recent models.
    \end{column}
    \begin{column}{0.5\textwidth}
      \begin{block}{Placement is not settled}
        \small
        \textbf{Gemma 3} uses four RMSNorms per block --- before \emph{and} after
        both sublayers. Arcee Trinity copied it.\\[0.4em]
        \textbf{Trinity} additionally initialises the second norm's gain at
        $\approx 1/\sqrt{L}$, so residual updates start small and grow.\\[0.4em]
        \textbf{Tiny Aya} (Cohere) deliberately \textbf{dropped} QK-Norm,
        because it can interact badly with long-context performance.
      \end{block}
    \end{column}
  \end{columns}
  \source{Raschka, \textit{A Dream of Spring for Open-Weight LLMs}, Feb 2026;
          \textit{Recent Developments in LLM Architectures}, May 2026.}
  \note{The Tiny Aya counter-example is the point of the slide. When we say
  ``standard'' here we only mean that most teams currently do it this way.
  Tell them to notice where a model report deviates, because that is usually
  where the team learned something.

  The formula, if you want it on the board: RMSNorm divides $x$ by
  $\sqrt{\tfrac{1}{d}\sum_i x_i^2 + \epsilon}$ and then multiplies
  element-wise by a learned gain vector $g$. The $\epsilon$ is there purely to
  stop division by zero. LayerNorm additionally subtracts the mean and adds a
  learned bias; dropping both is the entire difference, and it saves a pass
  over the vector.

  ``Gain initialised at $1/\sqrt{L}$'' in the Trinity line means residual
  updates start out small relative to the stream they are added to, and grow
  as training proceeds --- a stability trick, not an architectural one.}
\end{frame}

\begin{frame}[shrink=3]{Position: RoPE and its absence}
  Attention alone does not encode token order, so the model needs another
  source of positional information.
  \vspace{0.5em}
  \begin{columns}[T]
    \begin{column}{0.52\textwidth}
      \textbf{RoPE} --- rotate $q$ and $k$ by an angle proportional to position.
      The dot product then depends only on the \emph{relative} offset.
      \begin{itemize}
        \item no parameters
        \item extrapolates better than learned absolute embeddings
        \item extendable after training by scaling the base frequency (YaRN)
        \item the most common choice in current decoder models
      \end{itemize}
      \vspace{0.4em}
      \textbf{Partial RoPE} --- apply it to only a subset of dimensions
      (Qwen3-Next, Qwen3.5).
    \end{column}
    \begin{column}{0.48\textwidth}
      \begin{gotcha}[NoPE: sometimes none at all]
        Several 2026 models use \textbf{no positional encoding} in their
        \emph{global} attention layers: SmolLM3, Arcee Trinity, Kimi Linear.
        \vspace{0.4em}
        Why it works: another component already supplies position --- the
        sliding-window layers, or the linear-attention blocks, which are
        sequential by construction.
        \vspace{0.4em}
        Payoff: no RoPE retuning when you extend the context.
      \end{gotcha}
    \end{column}
  \end{columns}
  \note{NoPE is counter-intuitive enough to be worth the slide. The pattern to
  teach is conditional: you can drop positional encoding in a layer when
  something else in the model is already supplying the position information.
  No model drops it in every layer.

  YaRN reappears on Thursday under long-context mid-training.}
\end{frame}

\begin{frame}{The feed-forward network: SwiGLU}
  \[
    \mathrm{SwiGLU}(x) \;=\; \big(\mathrm{SiLU}(W_1 x) \odot W_3 x\big)\,W_2
  \]
  \begin{columns}[T]
    \begin{column}{0.54\textwidth}
      \begin{itemize}
        \item A \textbf{gated} activation: one branch decides \emph{what} to
          pass, the other decides \emph{how much}
        \item Three matrices instead of two, so the hidden width is conventionally
          cut to $\approx \tfrac{2}{3}$ to keep the parameter count equal
        \item The lineage: GeLU $\to$ SiLU $\to$ SwiGLU
        \item Where most of a dense model's parameters live
      \end{itemize}
    \end{column}
    \begin{column}{0.46\textwidth}
      \begin{block}{Why this matters later}
        \small
        The FFN is the big, parallel, per-position computation.
        \vspace{0.3em}
        That makes it the natural thing to \textbf{sparsify}, which is what
        Mixture of Experts does. We get there in ten minutes.
      \end{block}
    \end{column}
  \end{columns}
  \note{Do not spend long here. The two things that must land are the gating
  intuition and that the FFN is where the parameters are, because the second
  one sets up MoE.

  If someone asks why SwiGLU rather than something else: honestly, because it
  worked slightly better in Shazeer's ablations and everyone copied it. That
  is a true and useful answer about how this field operates.}
\end{frame}

\begin{frame}[shrink=20]{Four ways to reduce KV-cache costs}
  Many recent attention variants use one or more of these approaches.
  \vspace{0.6em}
  \centering
  \begin{tabular}{@{}L{2.7cm} L{3.4cm} L{5.6cm}@{}}
    \toprule
    \textbf{Move} & \textbf{Mechanism} & \textbf{Who does it}\\
    \midrule
    \key{Share heads} & MQA, GQA & Llama 3, Qwen3, Gemma 3, SmolLM3, MiniMax M2.5,
      Mistral Small\\[0.3em]
    \key{Compress the representation} & MLA, CCA & DeepSeek V3, GLM-5, Kimi K2,
      Mistral Large 3, Sarvam 105B, ZAYA1\\[0.3em]
    \key{Shorten or sparsify the sequence} & SWA, DSA, CSA + HCA & Gemma 3/4, OLMo 3,
      Arcee Trinity, DeepSeek V3.2 and V4\\[0.3em]
    \key{Share across layers} & cross-layer KV & Gemma 4 E2B / E4B\\
    \bottomrule
  \end{tabular}

  \vspace{0.8em}
  \raggedright
  These approaches can be combined. Gemma 4, for example, uses three of the
  four.
  \source{Raschka, \textit{A Visual Guide to Attention Variants}, Mar 2026;
          \textit{Recent Developments in LLM Architectures}, May 2026.}
  \note{Use the four rows as a taxonomy for the next five slides.

  The orthogonality claim is what earns the taxonomy its place. It means you
  can open a new model's config and decompose it into these moves instead of
  memorising a list of acronyms.}
\end{frame}

\begin{frame}{Share heads: MQA and GQA}
  \begin{columns}[T]
    \begin{column}{0.56\textwidth}
      Give several \emph{query} heads the same $K$ and $V$.
      \vspace{0.5em}

      \textbf{GQA is a spectrum:}
      \[
        \underbrace{h_{kv}=1}_{\text{MQA}} \;\longleftrightarrow\;
        \underbrace{h_{kv}=h_q}_{\text{MHA}}
      \]
      \begin{itemize}
        \item cache shrinks by $h_q/h_{kv}$
        \item fewer parameters too
        \item absolute memory savings grow with sequence length
        \item MQA saves the most cache, but may have the largest quality cost
      \end{itemize}
    \end{column}
    \begin{column}{0.44\textwidth}
      \begin{block}{Why it is still everywhere in 2026}
        \small
        Robust, simple to implement, few hyperparameters, and every inference
        stack has a fast kernel for it.
        \vspace{0.4em}
        MiniMax M2.5 and Nanbeige 4.1 use \textbf{only} GQA, with no other
        efficiency trick --- a deliberate 2026 choice.
      \end{block}
      \vspace{0.3em}
      \small
      Gemma 4 E2B goes all the way to MQA ($h_{kv}=1$).
    \end{column}
  \end{columns}
  \note{The spectrum framing is better than treating MQA, GQA and MHA as three
  separate things. It is one dial.

  The MiniMax M2.5 point is worth making: a 2026 flagship deliberately
  choosing the boring option is evidence that the exotic options are not
  free.}
\end{frame}

\begin{frame}{Compress the representation: MLA}
  \begin{columns}[T]
    \begin{column}{0.56\textwidth}
      MLA has the same goal as GQA but takes a different route: it stores a
      compressed representation of $K$ and $V$.
      \begin{itemize}
        \item project $K,V$ down into a low-dimensional \key{latent}
        \item cache the latent
        \item project back up at use time
      \end{itemize}
      \vspace{0.4em}
      Costs an extra matmul, saves memory and bandwidth.
      \vspace{0.6em}
      \textbf{Analogy:} the down- and up-projection is structurally LoRA.
      \vspace{0.5em}
      \textbf{Gotcha:} the \emph{query} is compressed too --- but only during
      training, not at inference.
    \end{column}
    \begin{column}{0.44\textwidth}
      \begin{block}{Evidence from DeepSeek-V2}
        \small
        In DeepSeek-V2's ablations, \textbf{GQA scores below MHA}, while
        \textbf{MLA matches or slightly beats MHA}.
        \vspace{0.4em}
        In that experiment, MLA reduced memory without the quality drop
        observed for GQA.
      \end{block}
      \begin{block}{A natural experiment}
        \small
        Sarvam shipped \textbf{30B with GQA} and \textbf{105B with MLA} from
        the same team on the same recipe, so MLA is an upgrade path you can
        actually take.
      \end{block}
    \end{column}
  \end{columns}
  \source{Raschka, \textit{A Visual Guide to Attention Variants}, Mar 2026;
          \textit{From DeepSeek V3 to V3.2}, Dec 2025.}
  \note{The DeepSeek-V2 ablation is, as far as anyone knows, still the most
  recent apples-to-apples comparison of these variants --- from 2024. Say that
  out loud, because it tells them something important about the state of the
  field: we are all making architecture decisions on two-year-old evidence.

  The Sarvam pair is the nicest piece of evidence in the lecture because it
  removes the ``maybe they just could not implement MLA'' confound.}
\end{frame}

\begin{frame}[shrink=9]{Shorten the sequence: sliding-window attention}
  Each token attends only to the last $t$ tokens. Interleave occasional
  \key{global} layers so information still crosses the whole sequence.
  \[
    O(n^2) \;\longrightarrow\; O(n \cdot t)
    \qquad\text{e.g. } n = 256{,}000,\ t = 4096
  \]
  \vspace{-0.4em}
  \centering
  \begin{tabular}{@{}l c c@{}}
    \toprule
    \textbf{Model} & \textbf{local : global} & \textbf{window}\\
    \midrule
    Gemma 2 & 1:1 & 4096\\
    Gemma 3 & 5:1 & 1024\\
    Gemma 4 E2B & 4:1 & ---\\
    OLMo 3, Arcee Trinity & 3:1 & 4096\\
    Xiaomi MiMo-V2-Flash & 5:1 & \textbf{128}\\
    Poolside Laguna XS.2 & 30 SWA : 10 global & 512\\
    \bottomrule
  \end{tabular}
  \vspace{0.7em}
  \begin{gotcha}
    ``Model X uses sliding-window attention'' is not enough to estimate its
    cost. You also need the window size and the ratio of local to global
    layers.
  \end{gotcha}
  \note{The Gemma 2 $\to$ Gemma 3 result is the interesting one. Everyone
  already knew local attention was cheaper. What surprised people was how far
  Google could push it, from a 1:1 ratio to 5:1 and a window from 4096 down to
  1024, while perplexity barely moved.

  SWA and GQA do different jobs and models use both. SWA reduces how much
  context a layer looks at. GQA reduces how much state each token contributes.
  Gemma 3 does both.}
\end{frame}

\begin{frame}{Sparsify the sequence: learned sparsity (DSA)}
  Same idea as a sliding window, except the subset is \key{learned} rather than
  fixed by locality.
  \vspace{0.5em}
  \begin{columns}[T]
    \begin{column}{0.54\textwidth}
      \textbf{DeepSeek Sparse Attention}, two stages:
      \begin{enumerate}
        \item \textbf{Lightning indexer} --- scores every past position for the
          new query, computed over MLA's already-compressed representations
        \item \textbf{Token selector} --- keeps the top-$k$, builds a sparse mask
      \end{enumerate}
      \vspace{0.4em}
      $k = 2048$ in the released code.
      \[
        O(L^2) \;\longrightarrow\; O(L\cdot k)
      \]
    \end{column}
    \begin{column}{0.46\textwidth}
      \begin{block}{MLA and DSA are orthogonal}
        \small
        MLA compresses \textbf{what is stored}.\\
        DSA reduces \textbf{how much is revisited}.\\
        DeepSeek V3.2 and GLM-5 use both.
      \end{block}
      \begin{block}{The stated goal}
        \small
        DeepSeek were not trying to beat V3.1-Terminus. They wanted the
        efficiency with as little \emph{degradation} as they could manage.
        Notice how modest that claim is.
      \end{block}
    \end{column}
  \end{columns}
  \source{Raschka, \textit{From DeepSeek V3 to V3.2}, Dec 2025.}
  \note{DSA is new and more complicated to implement than fixed local
  attention. In practice, kernel support often matters as much as an
  algorithm's asymptotic cost.

  DeepSeek V4 pushes this further with CSA and HCA, which compress \emph{along
  the sequence} so the cache itself gets shorter --- reportedly 10\% of V3.2's
  KV cache at 1M context. Mention it, do not teach it; flag that the paper has
  no ablation.}
\end{frame}

\begin{frame}{Share across layers, and replace attention entirely}
  \begin{columns}[T]
    \begin{column}{0.48\textwidth}
      \textbf{Cross-layer KV sharing} (Gemma 4)\\[0.3em]
      Later layers stop computing their own $K,V$ and reuse an earlier layer's
      --- but still compute their \emph{own queries}, so each layer keeps its
      own attention pattern.
      \vspace{0.4em}
      \small
      \begin{tabular}{@{}l c c@{}}
        \toprule
        & \textbf{own KV} & \textbf{saved @128k}\\
        \midrule
        E2B (35 layers) & first 15 & $\approx$ 2.7\,GB\\
        E4B (42 layers) & first 24 & $\approx$ 6\,GB\\
        \bottomrule
      \end{tabular}
      \vspace{0.4em}\\
      \textit{Cross-layer attention $\neq$ cross-attention.}
    \end{column}
    \begin{column}{0.52\textwidth}
      \textbf{Linear-attention hybrids} (Qwen3.5, Kimi Linear)\\[0.3em]
      Replace most full-attention layers with a recurrent, fixed-size state.
      Typically \key{3 linear : 1 full}.
      \[
        \text{state} = B \times h \times d_{\text{head}} \times d_{\text{head}} \times b
      \]
      \alert{No dependence on sequence length at all.}
      \vspace{0.4em}
      \small
      The price: a fixed-size state cannot do precise content-based retrieval,
      which is why one full-attention layer in four survives.
    \end{column}
  \end{columns}
  \note{Cross-layer sharing is an approximation --- it reduces model capacity.
  The cross-layer-attention paper reports minimal impact, but only for the
  small models they tested, and Gemma 4 E2B/E4B are small models. Be precise
  about that.

  For hybrids, teams lead with the constant-memory claim and tend to bury the
  retrieval weakness further down the report.}
\end{frame}

\begin{frame}{MiniMax: a change of direction}
  \centering
  \begin{tikzpicture}[node distance=9mm]
    \node[mod, minimum width=30mm] (m1) {MiniMax M1 \quad {\tiny Jun 2025}\\456B-A46B\\lightning attention};
    \node[mod, minimum width=30mm, right=of m1] (m2) {MiniMax M2\\230B\\\textbf{back to full attention}};
    \node[mod, minimum width=30mm, right=of m2] (m3) {MiniMax M2.5 \quad {\tiny Feb 2026}\\230B\\plain GQA only};
    \draw[flow] (m1) -- (m2);
    \draw[flow] (m2) -- (m3);
  \end{tikzpicture}
  \vspace{0.9em}

  \begin{columns}[T]
    \begin{column}{0.52\textwidth}
      The team reported that linear attention worked well on ordinary prompts
      but performed worse on \alert{reasoning and multi-turn tasks}. They
      therefore returned to full attention in M2.
    \end{column}
    \begin{column}{0.48\textwidth}
      The broader question remains open: Kimi Linear revisited the approach in
      late 2025, and
      \textbf{Qwen3.5 promoted the hybrid into its flagship line} in
      February 2026.
      \vspace{0.4em}
      Despite being ``only'' 230B, M2.5 has dominated OpenRouter usage.
    \end{column}
  \end{columns}
  \source{Raschka, \textit{Beyond Standard LLMs}, Nov 2025; \textit{A Dream of Spring}, Feb 2026.}
  \note{This is a useful example of architecture choices changing after
  deployment. Different teams can reach different conclusions because their
  workloads, implementations and evaluation criteria differ.

  It also sets up the honest closer on the next slide.}
\end{frame}

\begin{frame}{Comparing architectures carefully}
  \begin{gotcha}[There is no single ``best architecture'']
    Large, well-controlled public comparisons are rare. Most comparisons use
    \emph{trained models} that differ in data, optimisation, scale and
    architecture at the same time.
  \end{gotcha}
  \vspace{0.6em}
  \begin{takeaway}
    Architecture, data and training recipe all matter. Most of the changes in
    this section are motivated primarily by lower training or inference cost.
  \end{takeaway}
  \vspace{0.6em}
  \small
  A basic transformer block can fit in 50--100 lines of PyTorch, while these
  variants add substantial implementation complexity. A good way to learn
  them is to \textbf{start from GPT-2 and add one component at a time}, as you
  begin to do in Lab 1.
  \source{Raschka, \textit{A Visual Guide to Attention Variants}, Mar 2026;
          \textit{Recent Developments in LLM Architectures}, May 2026.}
  \note{The main point is that model comparisons confound many variables.
  Architecture papers are often motivated by cost, while data and training
  choices also have a large effect on quality. Thursday's lecture focuses on
  those choices.

  Clock check: should be around 17:15.}
\end{frame}

% =====================================================================
\section{Mixture of Experts}
% =====================================================================

\begin{frame}{Sparsify the FFN}
  Replace the single FFN with $N$ expert FFNs plus a small \key{router}.
  Each token is sent to the top-$k$ experts, usually $k \in \{1,\dots,8\}$.
  \vspace{0.5em}
  \begin{columns}[T]
    \begin{column}{0.5\textwidth}
      \begin{itemize}
        \item \key{Total parameters} --- what you must store and move
        \item \key{Active parameters} --- what actually runs per token, and what
          sets the FLOPs
        \item The model stores all experts, but computes with only the active
          subset for each token
      \end{itemize}
      \vspace{0.4em}
      \textbf{Shared experts} --- one or more experts that every token always
      uses, for the genuinely general computation. Common in 2026.
    \end{column}
    \begin{column}{0.5\textwidth}
      \begin{block}{Load balancing}
        \small
        Without some form of balancing, a few experts may receive most of the
        tokens while others are rarely used.
        \vspace{0.3em}
        An auxiliary loss, or a bias-based balancing scheme, keeps the
        distribution more even. This adds another source of complexity when
        training MoE models.
      \end{block}
    \end{column}
  \end{columns}
  \note{The key number is the ratio. DeepSeek V3 has 671B total parameters and
  37B active parameters,
  --- so you need the memory of a 671B model and the compute of a 37B model.
  This trade-off can work well in a data centre but is less attractive on a
  memory-constrained device.

  Load balancing is worth mentioning because it is a genuine reason small
  teams avoid MoE.}
\end{frame}

\begin{frame}[shrink=6]{The 2026 sparse models}
  \centering
  \begin{tabular}{@{}l r r l@{}}
    \toprule
    \textbf{Model} & \textbf{Total} & \textbf{Active} & \textbf{Notes}\\
    \midrule
    DeepSeek V3 / V3.2 & 671B & 37B & MLA; V3.2 adds DSA\\
    GLM-5 & 744B & 40B & 78 layers, 256 experts, MLA + DSA\\
    Qwen3.5 & 397B & 17B & hybrid attention, multimodal\\
    Qwen3-Next / Coder-Next & 80B & 3B & 262k native context\\
    Arcee Trinity Large & 400B & 13B & deliberately \emph{coarse} experts\\
    Kimi K2.5 & 1T & 32B & scaled-up DeepSeek V3 recipe\\
    Step 3.5 Flash & 196B & 11B & MTP-3 at inference\\
    MiniMax M2.5 & 230B & --- & plain GQA, no tricks\\
    \bottomrule
  \end{tabular}
  \vspace{0.8em}
  \begin{gotcha}
    Active parameter count is only part of the picture. Also check the KV
    cache, attention pattern, context length and routing overhead.
  \end{gotcha}
  \source{Raschka, \textit{A Dream of Spring}, Feb 2026; \textit{Recent Developments}, May 2026.}
  \note{Date this slide out loud. It is accurate as of the writing of this
  course and will be stale within a year, and that is a feature: part of the
  skill we are teaching is reading a new model report and slotting it into a
  framework.

  Note that none of these fit on one GH200, which is why their project models
  are 1--4B.}
\end{frame}

\begin{frame}{There is no agreed expert size}
  \begin{columns}[T]
    \begin{column}{0.48\textwidth}
      \begin{block}{Finer, more experts}
        \small
        GLM-4.7 $\to$ GLM-5: \textbf{160 $\to$ 256} experts\\[0.3em]
        Qwen3-Next: four times as many experts as Qwen3 235B, \emph{and} adds a
        shared expert --- while being 3$\times$ smaller\\[0.3em]
        ZAYA1-8B: routes to a \textbf{single} expert per token\\[0.3em]
        \textit{Rationale: more specialisation per unit of compute.}
      \end{block}
    \end{column}
    \begin{column}{0.48\textwidth}
      \begin{block}{Coarser, fewer experts}
        \small
        Arcee Trinity: started from a DeepSeek-style many-small-experts design
        and deliberately made it \textbf{coarser}\\[0.3em]
        Mistral 3 Large: the same move when adopting the DeepSeek V3
        architecture\\[0.3em]
        \textit{Rationale: inference throughput. Many tiny experts are a
        memory-access nightmare.}
      \end{block}
    \end{column}
  \end{columns}
  \vfill
  \begin{gotcha}[Relevant to your project choice]
    Gemma 4 ships a 31B \textbf{dense} model alongside its MoEs, explicitly
    ``for maximum quality and more convenient post-training, since MoEs are
    trickier to work with''.
    \vspace{0.3em}
    For this course project, prefer a dense model unless you have a specific
    reason to use MoE.
  \end{gotcha}
  \note{Different teams make opposite choices because they optimise for
  different constraints, such as model capacity or inference throughput.

  The gotcha is practical advice they should act on today.}
\end{frame}

\fi

% =====================================================================
\section{The model landscape}
% =====================================================================

\begin{frame}[shrink=3]{Useful distinctions between models}
  \centering
  \small
  \renewcommand{\arraystretch}{1.08}
  \begin{tabular}{@{}L{3.2cm} L{9.3cm}@{}}
    \toprule
    \textbf{Axis} & \textbf{What it means in practice}\\
    \midrule
    dense vs MoE & dense activates all parameters; MoE routes each token through a subset of expert layers\\
    \addlinespace[0.35em]
    open-weight vs open-source & releasing \emph{weights} does not necessarily
      mean releasing the data, code and full training recipe\\
    \addlinespace[0.35em]
    licence & Apache 2.0 / MIT are permissive; other licences may restrict particular uses\\
    \addlinespace[0.35em]
    base vs instruct & base continues text; instruct follows requests through a chat template. Week 2 compares both\\
    \addlinespace[0.35em]
    reasoning vs not & uses extra generation for intermediate reasoning; this
      can help on some difficult tasks, but adds latency and does not help uniformly\\
    \addlinespace[0.35em]
    decoder vs encoder & decoder models generate text; encoder models remain
      common for \emph{embeddings} and retrieval --- week 5\\
    \bottomrule
  \end{tabular}
  \note{Clarify the distinction between access to weights and access to the
  complete training pipeline. This matters both for reproducibility and for
  the licence discussion in the report.

  Flag the encoder line: after six weeks of decoders it surprises people that
  retrieval still runs on encoders.}
\end{frame}

\begin{frame}[shrink=3]{How to read a new model release}
  Use this checklist for your project and again in week 7.
  \vspace{0.5em}
  \begin{columns}[T]
    \begin{column}{0.5\textwidth}
      \begin{enumerate}
        \item Start with \texttt{config.json}, then read the release post
        \item Which variants exist: base, instruct, reasoning?
        \item How large is it, and what hardware does it require?
        \item What is its context length?
        \item Which tokenizer and chat template does it use?
      \end{enumerate}
    \end{column}
    \begin{column}{0.5\textwidth}
      \begin{enumerate}
        \setcounter{enumi}{5}
        \item How well does the tokenizer cover \emph{your} language?
        \item What does the licence allow?
        \item What data and post-training recipe are documented?
        \item Finally, inspect benchmarks and their evaluation setup
      \end{enumerate}
      \vspace{0.5em}
      \begin{gotcha}
        A model card is a starting point, not independent evidence. Test the
        model on data that represents your own task.
      \end{gotcha}
    \end{column}
  \end{columns}
  \note{The specific model names will date quickly, but this reading strategy
  should remain useful. Benchmark results make more sense once students know
  what model and evaluation setup produced them.

  Benchmark results come last deliberately. Week 4 explains how to evaluate
  models on a task of your own.}
\end{frame}

% =====================================================================
\section{Roihu in ten minutes}
% =====================================================================

\begin{frame}{Getting access}
  \begin{columns}[T]
    \begin{column}{0.52\textwidth}
      \begin{enumerate}
        \item Create a CSC account --- \textbf{multi-factor authentication is
          required} for the web interface
        \item Join the course project through the invitation link you were sent
        \item Accept the Roihu terms of use
        \item Open \texttt{www.roihu.csc.fi} and start the
          \textbf{Jupyter for courses} app
      \end{enumerate}
      \vspace{0.5em}
      The app clones the course repository and loads the course environment for
      you. Nothing to install.
    \end{column}
    \begin{column}{0.48\textwidth}
      \begin{block}{What you get}
        \small
        \textbf{Roihu} --- CSC's national supercomputer, in service since
        summer 2026 (Puhti and Mahti are retired).\\[0.4em]
        GPU partition: 132 nodes, four NVIDIA \textbf{GH200} superchips each.\\[0.4em]
        One GH200 = one H100-class GPU with \textbf{96\,GiB} + a 72-core
        \textbf{ARM Grace} CPU.
      \end{block}
    \end{column}
  \end{columns}
  \vspace{0.4em}
  \alert{Please do this today.} Account approval can take time, and Lab 1 uses
  Roihu.
  \note{Account approval may take time, so students should register today.

  If there is time, do this live on the projector.}
\end{frame}

\begin{frame}[shrink=9]{Two ways to run}
  \centering
  \begin{tabular}{@{}L{3.1cm} l L{3.3cm} L{3.7cm}@{}}
    \toprule
    \textbf{Mode} & \textbf{Partition} & \textbf{Limits} & \textbf{Used for}\\
    \midrule
    Interactive Jupyter, CPU & \texttt{interactive} & 36 h, up to 32 cores &
      tokenizers, data prep, writing code\\[0.3em]
    Interactive Jupyter, GPU & \texttt{gpuinteractive} & 12 h, one GH200 &
      lab sessions, prompting, LoRA SFT, vLLM\\[0.3em]
    Batch, short & \texttt{gputest} & 15 min & \textbf{debugging your script}\\[0.3em]
    Batch, long & \texttt{gpumedium} & 36 h, 1--4 GPUs &
      tiny-GPT pretraining, GRPO, long fine-tunes\\
    \bottomrule
  \end{tabular}
  \vspace{0.9em}
  \begin{gotcha}[The GPU nodes use ARM CPUs]
    The GPU nodes run \textbf{ARM}. Build any virtual environment on Roihu's
    ARM login nodes, and check that any additional package supports the
    platform before depending on it.
  \end{gotcha}
  \note{Repeat the practical rule: test on \texttt{gputest} before submitting
  to \texttt{gpumedium}. This catches simple errors with a much shorter wait.

  The ARM point will come back in weeks 2 and 3. Say it now so it is not a
  surprise.}
\end{frame}

\begin{frame}[fragile]{A batch job, line by line}
\begin{lstlisting}[language=slurm]
#!/bin/bash
#SBATCH --account=project_2021893   # the course project
#SBATCH --partition=gputest         # <- gputest first, always
#SBATCH --nodes=1
#SBATCH --gres=gpu:1                # one GH200
#SBATCH --cpus-per-task=16
#SBATCH --mem=64G
#SBATCH --time=00:15:00             # gputest caps at 15 minutes
#SBATCH --output=logs/%x-%j.out     # job name and job id in the filename

module purge
module load pytorch                 # the course environment

srun python train_tiny_gpt.py --config configs/small.yaml
\end{lstlisting}
  \vspace{0.3em}
  \small
  \texttt{sbatch job.sh} to submit \quad
  \texttt{squeue --me} to watch \quad
  \texttt{scancel <jobid>} to stop \quad
  \texttt{seff <jobid>} to see what you actually used
  \note{Walk every line. Most of the room has never written a Slurm script and
  the ones who have will forgive ninety seconds.

  \texttt{seff} is the one people do not know and should: it tells them
  whether they asked for four times the memory they used, which matters when
  twenty-five people share an allocation.}
\end{frame}

\begin{frame}[shrink=15]{Lab 1 --- take-home work}
  \begin{columns}[T]
    \begin{column}{0.56\textwidth}
      \textbf{A. Tokenizer} \hfill \textit{CPU Jupyter}\\
      Train a byte-level BPE on Finnish Wikipedia at 2--3 vocabulary sizes,
      then compare fertility against Qwen3.5 and Gemma 4.
      \vspace{0.25em}

      \textbf{B. Tiny GPT} \hfill \textit{GPU batch}\\
      Train a 10--30M-parameter decoder with your tokenizer. Script and
      Slurm template provided; about 20 min per run.
      \vspace{0.25em}

      \textbf{C. Scaling}\\
      Three model sizes at a fixed token budget. Plot validation loss
      against compute.
      \vspace{0.25em}

      \textbf{D. What scale buys}\\
      Run the same few-shot prompt on your tiny model and on a larger
      pretrained base model.
      \vspace{0.25em}

      \textbf{E. Project}\\
      Choose a task, model and held-out set.
    \end{column}
    \begin{column}{0.44\textwidth}
      \begin{block}{In the final report}
        \small
        a fertility table \textbullet\ loss curves \textbullet\ a scaling plot
        \textbullet\ one paragraph comparing your run's compute with a published
        frontier run
      \end{block}
      \begin{gotcha}[There is no lab session this week]
        Thursday is a lecture, so start early and use the course channel.
      \end{gotcha}
      \vspace{0.2em}
      \small
      The scaling plot is explained on Thursday. You can start A, B and E
      tonight.
    \end{column}
  \end{columns}
  \note{Be explicit that Thursday is a lecture, not a lab. Every year someone
  turns up on Thursday expecting help with the notebook.

  Tell them the order: A and E tonight, B once Thursday's lecture has
  explained the training loop, C and D at the weekend.}
\end{frame}

% =====================================================================
\begin{frame}[standout]
  Three things to remember
\end{frame}

\begin{frame}{Three things to remember}
  \begin{enumerate}
    \item \key{Language models predict distributions over the next token.}
      Generating text means choosing a token, adding it to the context and
      repeating.
      \vspace{0.6em}
    \item \key{Tokenization and architecture shape both quality and cost.}
      Transformers build contextual token representations with attention and
      feed-forward layers.
      \vspace{0.6em}
    \item \key{An LLM application is more than model weights.} Post-training,
      prompting, retrieval, tools and evaluation all affect what users see.
  \end{enumerate}
  \vfill
  \textbf{Thursday:} where the weights come from. Data pipelines, scaling laws,
  and what a real pretraining run looks like.
  \vspace{0.3em}

  \textbf{Tonight:} create your CSC account and start choosing your task and
  model.
  \note{These three sentences get repeated in the week 7 wrap-up, so say them
  cleanly.

  End by reminding students to create their CSC account today.}
\end{frame}

\appendix

\begin{frame}{Reading}
  \begin{columns}[T]
    \begin{column}{0.5\textwidth}
      \textbf{Core}
      \begin{itemize}\small
        \item Stanford CS336, \textit{Language Modeling from Scratch}
        \item Karpathy, \textit{nanoGPT} and \textit{nanochat}
        \item Vaswani et al. 2017, \textit{Attention Is All You Need}
      \end{itemize}
      \vspace{0.5em}
      \textbf{Architecture, current}
      \begin{itemize}\small
        \item Raschka, \textit{A Visual Guide to Attention Variants in Modern LLMs} (Mar 2026)
        \item Raschka, \textit{Recent Developments in LLM Architectures} (May 2026)
        \item Raschka, \textit{A Dream of Spring for Open-Weight LLMs} (Feb 2026)
      \end{itemize}
    \end{column}
    \begin{column}{0.5\textwidth}
      \textbf{Specific mechanisms}
      \begin{itemize}\small
        \item Ainslie et al. 2023 --- GQA
        \item DeepSeek-V2 / V3 technical reports --- MLA, MoE
        \item Raschka, \textit{From DeepSeek V3 to V3.2} (Dec 2025) --- DSA
        \item Su et al. 2021 --- RoPE
        \item Shazeer 2020 --- GLU variants
      \end{itemize}
      \vspace{0.5em}
      \textbf{Roihu}
      \begin{itemize}\small
        \item \texttt{docs.csc.fi/computing/systems-roihu}
        \item \texttt{docs.csc.fi/computing/webinterface/}\\
              \texttt{jupyter-for-courses}
      \end{itemize}
    \end{column}
  \end{columns}
  \note{Point at CS336 as the thing to follow if they want the full
  from-scratch treatment, and at Raschka's architecture posts as the way to
  stay current after the course ends.}
\end{frame}

\end{document}
